\section*{Capítulo \ref{ch:b1:active-margins} --- Márgenes activos}

\begin{solution}{exo:b1:active-margins:1}
(a) $2\pi G\,\Delta\rho\,h = 2\pi\times\num{6.674e-11}\times1700\times3000 =
\qty{2.14e-3}{m.s^{-2}}$: la gravedad sería menor en unos
\qty{214}{\milli\gal}. (b) Para \qty{3.5}{km},
$2\pi\times\num{6.674e-11}\times1700\times3500 = \qty{2.50e-3}{m.s^{-2}}$,
\qty{250}{\milli\gal}: el tamaño de la depresión observada, así que la fosa está lejos del equilibrio isostático. (c) La fórmula de la losa supone que la roca que falta
se extiende sin límite en todas direcciones; una fosa de solo decenas de kilómetros de anchura quita mucha menos masa alrededor del punto de medida, así que
su agua sola da menos de \qty{250}{\milli\gal}. El resto del déficit yace más abajo: toda la placa se curva hacia abajo bajo la fosa, y su \omterm{def:b1:layered-earth:layers}{corteza}, más
ligera que el \omterm{def:b1:layered-earth:layers}{manto}, ocupa el lugar del \omterm{def:b1:layered-earth:layers}{manto} en su base.
\end{solution}

\begin{solution}{exo:b1:active-margins:2}
(a) $\qty{82}{mm.yr^{-1}} = \qty{82}{km.Myr^{-1}}$, así que \qty{820}{km} en
\qty{10}{Myr}. (b) $820\sin\ang{45} = \qty{580}{km}$. (c) El camino hasta
\qty{600}{km} es $600/\sin\ang{45} = \qty{849}{km}$, recorrido en
$849/82 = \qty{10.3}{Myr}$.
\end{solution}

\begin{solution}{exo:b1:active-margins:3}
$g = 110/\tan\delta$: \qty{191}{km} para \ang{30}, \qty{110}{km} para \ang{45},
\qty{64}{km} para \ang{60}. El arco sobre la losa de \omterm{def:b1:geological-maps:attitude}{buzamiento} más suave yace más lejos
de su fosa.
\end{solution}

\begin{solution}{exo:b1:active-margins:4}
(a) Convergencia: $\qty{2.7e6}{m^2.yr^{-1}}/\qty{5.1e7}{m} =
\qty{0.053}{m.yr^{-1}}$, unos \qty{53}{mm.yr^{-1}}. Expansión:
$\qty{3.0e6}{m^2.yr^{-1}}/\qty{6.7e7}{m} = \qty{0.045}{m.yr^{-1}}$, unos
\qty{45}{mm.yr^{-1}}. Con una longitud total menor, las fronteras de subducción consumen fondo oceánico más deprisa de lo que las dorsales lo producen, pero consumen menos en total.
(b) El \qty{0.3}{km^2} restante aproximadamente se absorbe a lo largo de las otras fronteras convergentes del modelo, como las colisiones continentales y las zonas convergentes dentro de los océanos, donde las placas se acortan en lugar de subducirse; y las dorsales no son las únicas fronteras divergentes: los \omterm{def:b1:rifts-ridges-margins:rift}{rifts continentales} y algunas transformantes también se abren un poco.
\end{solution}

\begin{solution}{exo:b1:active-margins:5}
(a) $\bar x = 350$, $\bar z = 290$; las desviaciones de $x$ son $-150$, $-50$,
$50$, $150$, y las de $z$ son $-140$, $-40$, $40$, $140$, así que
$S_{xx} = 2(150^2 + 50^2) = \num{50000}$ y
$S_{xz} = 2(150\times140 + 50\times40) = \num{46000}$. Entonces $s = 0.92$,
$a = 290 - 0.92\times350 = \qty{-32}{km}$, y el \omterm{def:b1:geological-maps:attitude}{buzamiento} es
$\arctan 0.92 = \ang{42.6}$. (b) $z = 0$ en $x = 32/0.92 = \qty{34.8}{km}$
de la fosa. (c) Con los puntos extremos la pendiente es $280/300 = 0.933$, un
\omterm{def:b1:geological-maps:attitude}{buzamiento} de \ang{43.0}: cercano, pero usa solo dos de los cuatro terremotos.
\end{solution}

\begin{solution}{exo:b1:active-margins:6}
(a) $50/\tan\ang{25} + 60/\tan\ang{45} = 107.2 + 60.0 = \qty{167.2}{km}$.
(b) Una losa de \ang{45} en todo su recorrido da \qty{110}{km}: el segmento somero y suave ensancha la brecha en
\qty{57}{km}. (c) El segmento profundo es $z = 50 + (x - 107.2)$; alcanza $z = 0$ en $x = \qty{57.2}{km}$, hacia tierra desde
la fosa, como hace la recta de las Kuriles a \qty{53}{km}: un segmento somero suave explica el desplazamiento de las Kuriles. (d) $107.2 + 60/\tan\ang{60} = 107.2 + 34.6 =
\qty{141.9}{km}$.
\end{solution}

\begin{solution}{exo:b1:active-margins:7}
(a) El radio es la mitad del \omterm{def:b1:geological-maps:attitude}{buzamiento}, y $r = R\,\theta$ con $\theta$ en radianes:
\ang{22.5} y $6371\times0.3927 = \qty{2502}{km}$ para \ang{45}; \ang{15} y
\qty{1668}{km} para \ang{30}; \ang{35} y \qty{3892}{km} para \ang{70}.
(b) Los arcos sobre losas de \omterm{def:b1:geological-maps:attitude}{buzamiento} más suave. (c) La predicción es
$44.9/2 = \ang{22.45}$; el ajuste, \ang{21.9}, es
$0.55/22.45 = \qty{2.4}{\percent}$ más pequeño.
\end{solution}

\begin{solution}{exo:b1:active-margins:8}
(a) $q/k$: $0.035/2.5 = \qty{0.014}{K.m^{-1}} = \qty{14}{K.km^{-1}}$ bajo
el \omterm{def:b1:active-margins:active-margin}{antearco}, $0.080/2.5 = \qty{32}{K.km^{-1}}$ bajo el tras-arco.
(b) $10 + 14\times30 = \qty{430}{\celsius}$ y
$10 + 32\times30 = \qty{970}{\celsius}$. (c) Parte del \omterm{def:b1:heat-budget:heat-flow}{flujo de calor} en superficie
lo produce la radiactividad dentro de la propia \omterm{def:b1:layered-earth:layers}{corteza}, así que menos calor atraviesa
los niveles más profundos y el gradiente disminuye con la profundidad: una extrapolación lineal del
gradiente en superficie sobrestima la temperatura en profundidad.
\end{solution}

\begin{solution}{exo:b1:active-margins:9}
(a) \qty{32}{mm.yr^{-1}} es \qty{32}{km.Myr^{-1}}; la velocidad media es
\qty{16}{mm.yr^{-1}}. (b) $100/32 = \qty{3.1}{Myr}$. (c) La velocidad no tiene por qué
haber sido constante; parte de la anchura puede venir del estiramiento y el rifting del arco antes de que empezara la expansión;
y el eje puede haber saltado o haberse expandido de forma asimétrica.
\end{solution}

\begin{solution}{exo:b1:active-margins:10}
(a) El camino es $600/\sin\ang{45} = \qty{849}{km}$, recorrido a
\qty{82}{km.Myr^{-1}} en \qty{10.3}{Myr}, es decir \qty{3.27e14}{s}.
(b) $\sqrt{\num{e-6}\times\num{3.27e14}} = \qty{1.8e4}{m}$, unos
\qty{18}{km}. (c) $(\qty{5e4}{m})^2/\qty{e-6}{m^2.s^{-1}} = \qty{2.5e15}{s}$,
unos \qty{79}{Myr}. (d) El calor solo ha llegado a los \qty{18}{km} exteriores de cada
cara de la losa: su núcleo llega a los \qty{600}{km} casi tan frío como cuando salió
de la fosa, más frío y más resistente que el \omterm{def:b1:layered-earth:layers}{manto} que lo rodea, lo que explica por qué
aún puede romperse en terremotos a esa profundidad.
\end{solution}

\begin{solution}{exo:b1:active-margins:11}
(a) A la distancia $L$ detrás del pie, la superficie se alza $L\tan\alpha$ sobre el nivel del pie
y la base yace $L\tan\beta$ por debajo: la parte posterior del prisma tiene una altura de
$L(\tan\alpha + \tan\beta)$, y el triángulo tiene el área
$\frac12 L\times L(\tan\alpha + \tan\beta)$. (b) $\tan\ang{4} + \tan\ang{6} =
0.175$. El prisma recibe $\qty{1}{km}\times\qty{40}{km.Myr^{-1}} =
\qty{40}{km^2.Myr^{-1}}$ de sección, así que $A = 40t$ y
$L = \sqrt{2A/0.175}$: \qty{21.4}{km} tras \qty{1}{Myr}, \qty{42.8}{km}
tras \qty{4}{Myr}, \qty{67.6}{km} tras \qty{10}{Myr}. $L$ crece como
$\sqrt t$: cuatro veces más tiempo da el doble de anchura. (c) $A$ se reduce a la mitad y $L$
se divide por $\sqrt2$: \qty{15.1}{km}, \qty{30.2}{km} y \qty{47.8}{km}.
\end{solution}

\begin{solution}{exo:b1:active-margins:12}
(a) $\partial F/\partial a = -2\sum(z_i - a - s x_i) = 0$ da
$n\bar z = na + sn\bar x$, es decir $\bar z = a + s\bar x$. (b) Cada $z_i$ y
$\bar z$ crecen \qty{10}{km}, así que las desviaciones $z_i - \bar z$ y $S_{xz}$
no cambian: la pendiente es la misma, y $a = \bar z - s\bar x$ crece en
\qty{10}{km}. (c) Las desviaciones $x_i - \bar x$ no cambian, así que la pendiente es
la misma; $\bar x$ disminuye en 20, así que $a$ aumenta en $20s$. (d) Ninguno:
el \omterm{def:b1:geological-maps:attitude}{buzamiento} depende solo de la pendiente. Los errores compartidos por todos los puntos
mueven la recta sin girarla.
\end{solution}

\begin{solution}{pb:b1:active-margins:1}
\textbf{1.} Los ocho puntos yacen cerca de una recta que buza unos
\ang{45}, desde \qty{96}{km} de profundidad a \qty{155}{km} de la fosa hasta
\qty{573}{km} de profundidad a \qty{609}{km}.

\textbf{2.} $\sum x_i = 3152$, $\sum z_i = 2751$, $\sum x_i^2 = \num{1442910}$,
$\sum x_i z_i = \num{1290548}$; $\bar x = \qty{394.0}{km}$,
$\bar z = \qty{343.875}{km}$.

\textbf{3.} $S_{xx} = 1442910 - 8\times394^2 = \num{201022}$ y
$S_{xz} = 1290548 - 8\times394\times343.875 = \num{206654}$. Entonces
$s = 206654/201022 = 1.028$ y $a = 343.875 - 1.02802\times394 =
\qty{-61.2}{km}$: $z = -61.2 + 1.028\,x$.

\textbf{4.} $\arctan1.028 = \ang{45.8}$.

\textbf{5.} En $x = 61.2/1.028 = \qty{59.5}{km}$ hacia tierra desde la fosa. Una losa que hubiera salido de la fosa con este \omterm{def:b1:geological-maps:attitude}{buzamiento}
llegaría a la superficie allí; llega a ella en la fosa, así que buza más suavemente cerca de la fosa y se empina debajo.

\textbf{6.} Los desajustes son $-2.1$, $2.0$, $-7.3$, $5.3$, $7.1$, $2.3$,
$-15.0$ y $8.1$ \unit{km}. El terremoto 7 yace \qty{15}{km} por encima de la recta, el 8
\qty{8}{km} por debajo; los demás yacen a unos \qty{7}{km} de ella, una fracción
del espesor de la losa.

\textbf{7.} $\ang{320.9} - \ang{300} = \ang{20.9}$. Perpendicular a la
fosa: $87.4\cos\ang{20.9} = \qty{81.6}{mm.yr^{-1}}$; paralela:
$87.4\sin\ang{20.9} = \qty{31.2}{mm.yr^{-1}}$. La convergencia es oblicua.

\textbf{8.} \qty{81.6}{mm.yr^{-1}} es \qty{81.6}{km.Myr^{-1}}: \qty{816}{km} en
\qty{10}{Myr}.

\textbf{9.} El camino es $632/\sin\ang{45} = \qty{894}{km}$, recorrido en
$894/81.6 = \qty{11.0}{Myr}$.

\textbf{10.} $122 + 11 = \qty{133}{Myr}$.

\textbf{11.} $t = \qty{10.95}{Myr} = \qty{3.46e14}{s}$ y
$\sqrt{\num{e-6}\times\num{3.46e14}} = \qty{1.86e4}{m}$, unos \qty{19}{km}
de cada cara: mucho menos que los \qty{50}{km} hasta el centro de una losa de
\qty{100}{km} de espesor, cuyo núcleo sigue frío.

\textbf{12.} $q/k$: \qty{17.2}{K.km^{-1}} bajo el fondo oceánico antiguo,
\qty{14.0}{K.km^{-1}} bajo el \omterm{def:b1:active-margins:active-margin}{antearco} y \qty{32.0}{K.km^{-1}} bajo el
tras-arco.

\textbf{13.} Bajo el \omterm{def:b1:active-margins:active-margin}{antearco} la losa fría yace solo a unas decenas de kilómetros de profundidad; absorbe calor de la roca que tiene encima,
así que llega menos calor a la superficie.

\textbf{14.} $h = \qty{2.60e-3}{m.s^{-2}}/(2\pi\times\num{6.674e-11}\times1700)
= \qty{3647}{m}$, unos \qty{3.6}{km}: tanto como, o un poco más que, toda la profundidad de la fosa bajo el fondo que la rodea, unos
\qty{3.5}{km}. La fosa carece de masa que nada compensa.

\textbf{15.} La diferencia de densidad es $3300 - 1000 = \qty{2300}{kg.m^{-3}}$,
y $2\pi\times\num{6.674e-11}\times2300\times3500 = \qty{3.38e-3}{m.s^{-2}}$,
\qty{338}{\milli\gal}. La fórmula de la losa cuenta la masa que falta hasta cualquier
distancia; la fosa real está curvada hacia abajo solo en una franja de decenas de kilómetros
de anchura, así que su déficit es menor, unos \qty{260}{\milli\gal}, entre las dos
estimaciones de la losa.

\textbf{16.} La resistencia de la placa curvada, que mantiene la fosa hacia abajo y la elevación exterior hacia arriba como lo haría una placa
elástica, y el tirón de la losa densa que se hunde en el \omterm{def:b1:layered-earth:layers}{manto}.

\textbf{17.} $-61.2 + 1.028\times164 = \qty{107.4}{km}$.

\textbf{18.} Sí: \qty{107}{km} yace entre 72 y \qty{173}{km}, y solo unos
\qty{2}{km} más profundo que el promedio, \qty{105}{km}.

\textbf{19.} $x = (100 + 61.2)/1.028 = \qty{156.8}{km}$.

\textbf{20.} $\ang{45.8}/2 = \ang{22.9}$, y
$6371\times22.9\times\pi/180 = \qty{2546}{km}$.

\textbf{21.} El ajuste, \ang{21.9}, es \ang{1.0} menor, unos
\qty{4}{\percent}: la cáscara curvada describe bien la curvatura del arco de las Kuriles.

\textbf{22.} $x = (105 + 61.2)/1.028 = 166.2/1.028 = \qty{161.7}{km}$.

\textbf{23.} $164 - 161.7 = \qty{2.3}{km}$ más cerca de la fosa que el frente
observado, $2.3/164 = \qty{1.4}{\percent}$: los volcanes de las Kuriles se alzan sobre una losa
algo más profunda que el promedio, como encontró la pregunta 18.

\textbf{24.} Al kilómetro más próximo, la recta de la losa de las Kuriles predice una \omterm{def:b1:active-margins:volcanic-arc}{brecha arco--fosa} de
\textbf{\qty[text-series-to-math=true]{162}{km}} para volcanes sobre la losa a \qty{105}{km} (observada: \qty{164}{km}).
\end{solution}
